\documentclass[journal,compsoc,10pt]{IEEEtran}

% Load some possibly useful packages. You can remove the ones you don't need or
% add other packages that you need in your paper.
\usepackage{graphicx}                   % required to include graphics
\usepackage{amsfonts,amssymb,amsmath}   % mathsymbols etc
\usepackage{mathtools}                  % includes some fixes to amsmath, and adds convenience macros
\usepackage{bm}                         % make any letter boldface in math
\usepackage{algpseudocode}              % aka `algorithmicx`: same look `algorithmic`, but more flexible
\usepackage{array}                      % for better looking arrays
\usepackage{booktabs}                   % for better looking tables
\usepackage[usenames]{xcolor}           % colors for highlighting etc.
\usepackage{url}                        % for urls
\usepackage[%pdftitle=,                 % define the PDF metadata
            %pdfauthor=,
            colorlinks=true, linkcolor=purple, urlcolor=blue, citecolor=cyan,
            anchorcolor=black   % make links/references 'clickable' without ugly frames
            ]{hyperref}            
\usepackage[capitalise]{cleveref}                   % smart references to figures, equations etc.

% some packages which may be useful when drafting, but should be removed before submission
\usepackage{todonotes}

%own packages
\usepackage{amsthm}
\usepackage{float}
%\usepackage{subcaption}
\usepackage{stmaryrd}

% You can define your own commands. This is useful for shorthands, for example
\newcommand{\Exp}{\mathbb{E}}
\newcommand{\RR}{\mathbb{R}}
\newcommand{\CC}{\mathbb{C}}
\newcommand{\NN}{\mathbb{N}}
\newcommand{\MC}{\mathcal{M}} %Markov Chain
\newcommand{\IMC}{\mathcal{I}} %Interval MC
\newcommand{\tT}{\mathcal{T}} %Transition Matrix
\newcommand{\eE}{\mathcal{E}} 
\newcommand{\gG}{\mathcal{G}} %graph
\newcommand{\gS}{\mathcal{S}} %sinks
\newcommand{\nN}{\mathcal{N}}
\newcommand{\PF}{\mathbf{P}} % Probability function
\newcommand{\PFupper}{\mathbf{\overline{P}}}
\newcommand{\PFlower}{\mathbf{\underline{P}}}
\newcommand{\Initupper}{\mathbf{\overline{\iota}}_{init}}
\newcommand{\Initlower}{\mathbf{\underline{\iota}}_{init}}
\newcommand{\cCircle}{\bigcirc} % Probability function
\newcommand{\Path}{Paths^{\MC}_{inf}} % Probability function
\newcommand{\lnorm}[1]{\ell_{#1}\text{-norm}}

% Commands are also useful to ensure standardized notation, such as how to display vectors and matrices.
% The point is to replace formatting commands with semantic ones.
% This both makes the source more semantic, and also makes it easier to change formatting choices later
% IMPORTANT: In TeX, all commands share a global namespace. You need to be careful therefore not to overwrite an existing command.
%            Use \newcommand to ensure you don't accidentally overwrite an important macro
%            Use \renewcommand if you *do* want to overwrite a macro, have checked what the original does,
%            and have made sure the redefinition is safe.
\newcommand{\T}{\mathsf{T}}            % transpose
\renewcommand{\vec}[1]{\bm{#1}}        % redefine \vec to produce boldface instead of adding an arrow
\newcommand{\mat}[1]{\bm{#1}}          % define \vec to produce boldface
\newcommand{\ones}{\bm{1}}             % all-ones vector
\newcommand{\zeros}{\bm{0}}             % all-ones vector
\newcommand{\zerosmatr}{\bm{0_{n\times n}}}
\newcommand{\onessmatr}{\bm{1_{n\times n}}}     % all-zeros matrix
\newcommand{\eye}{\bm{I}}              % identity matrix
\newcommand{\ei}{\bm{e}_i}             % standard basis vector e_i
\newcommand{\ej}{\bm{e}_j}             % standard basis vector e_j
    % TIP: Resist the temptation to be too clever with macros
    %   - \ei and \ej are easier to type than e.g. \ee[i]  (note that \e is a low-level command and should not be redefined)
    %   - While you can define macros in terms of other macros  (eg we could define \ei in terms of \vec),
    %     many journals won’t accept this. Recursive macros also produce more frequent and more complicated compilation errors

% eigene libs
\newtheorem{theorem}{Theorem}[section]
\newtheorem{corollary}{Corollary}[theorem]
\newtheorem{lemma}{Lemma}[section]
\newtheorem{definition}{Definition}[section]
\usepackage{enumitem}

% Commands like \sin and \exp ensure proper fonts and spacing for math operators.
% Use the dedicated \DeclareMathOperator to define new operators
\DeclareMathOperator{\diag}{diag}

\begin{document}%[tikz,border=3mm]{standalone}
%\usetikzlibrary{arrows.meta,positioning,quotes}

\title{Expected Visiting Times of Interval Markov Chains}

% author information
% COMMENT OUT THESE LINES FOR YOUR CONFERENCE SUBMISSION!
\author{Jesper Nilsson, 457057}

\maketitle
\thispagestyle{plain}
\pagestyle{plain}

\begin{abstract}
    test
\end{abstract}


\section{Introduction}
Thema einführen 

\section{Preliminaries}
\subsection{Notation}\label{sec:notation}
We denote the entry in the $i$-th row and $j$-th column of a matrix $\tT$, as $(\tT)_{ij}$. We will use the notation $\underline{n} = {1, 2, \dots, n}$ with $n\in \NN$ for sets.
Cardinality of set $S$ as $|S|$.

$\cCircle$
$\Diamond$

\subsection{Matheteil}
Linear Order
Measurable set (used in \cref{def:next_step_probability})
Iverson Brackets
%Matrix Theory
\begin{definition}[Row-Stochastic Matrix]\label{def:row-stoch_matrix}
    A square matrix $\tT \in \RR^{n\times n}$ is row-stochastic, if and only if $\sum_{j=0}^n(\tT)_{ij} = 1$ for all rows $i \in \underline{n}$
\end{definition}
When considering each row as a function $f_i(j)=(\tT)_{ij}$, these functions are discrete probability distributions.
%Probability Theory
Expected Values
\begin{definition}[Discrete Probability Distribution]
    A function $P$ is a discrete probability distribution, if $P: X \rightarrow [0,1]$ where X is a countable set and $\sum_{x\in X}P(x) = 1$.
\end{definition}


\subsection{Markov Chains}
OMISSION OF $^\MC$

In this Section, we provide definitions and ... for Markov Chains, as well as Interval Markov Chains. The notation is based on \cite{Bibel}, \cite{Semantics_Model-Checking_omega-Regular}...
\subsubsection{Discrete-Time (Interval) Markov Chain}
\begin{definition}[DTMC]\label{def:DTMCs}
    A discrete-time Markov chain (DTMC) is a tuple $\MC = (S, \PF,\iota_{init})$, where
    \begin{enumerate}[label=(\roman*)]
        \item $S$ is a countable, nonempty set of states,
        \item $\PF: S \times S \rightarrow [0,1]$ is the transition probability function, such that $\sum_{s'\in S}\PF(s,s') = 1$ for all states $s \in S$, and
        \item $\iota_{init}: S \rightarrow [0,1]$ is the initial distribution, such that $\sum_{s\in S}\iota_{init}(s)=1$.
    \end{enumerate}
\end{definition}
We assume in this thesis, that the DTMCs are finite, i.e., that the set of states $S$ is finite. Furthermore, we assume an arbitrary but fixed linear order $\prec$ on the set $S$, such that the states $s\in S$ can be indexed. Thus, $s_i \prec s_j$ if $s,s' \in S$ and $i < j$.
Due to the conditions (i) and (ii) provided in \cref{def:DTMCs}, the transition probability function $\PF$ and the initial distribution $\iota_{init}$ are discrete probability distributions.

??(notion of a current state)
Moreover, if $\PF(s,s')>0$ for two given $s,s' \in S$, then it is possible to move from $s$ to $s'$ in one step. We refer to these states $s'$ as direct successors of $s$ and to $s$ as a direct predecessor of all states $s'$. Additionally, we define a neighbour set $N$, with $s' \in N(s)$, if and only if 

PROBABILITY SPACE OF DTMC

\begin{definition}[Transition Matrix]\label{def:transition_matrix}
    The transition matrix $\tT^{\MC}$ of a DTMC $\MC$ is a square matrix with size $S \times S$ \todo{gemeintes S klar?}, such that $(\tT^{\MC})_{ij} = \PF(s_i,s_j)$.
\end{definition}
Notice that due to the condition, given in (ii) of \cref{def:DTMCs}, that $\sum_{s_j\in S}\PF(s_i,s_j) = 1$ for all states $s \in S$, by defintion $\sum_{j=0}^{|S|}(\tT)_{ij} = 1$ for all rows $i \in \underline{n}$. Thus, the transition matrix $\tT$ of a DTMC $\MC$ is row-stochastic, see \cref{def:row-stoch_matrix}.

\begin{definition}
    Every DTMC $\MC$ induces a directed weigthed graph $\gG^{\MC} = (V,E,f)$, with the nodes $V = S^{\MC}$, the edges $(s,s')\in E$ if and only if $\PF(s,s')>0$, and the weight function $f(s,s') = \PF(s,s')$\todo{simply use $\PF$ instead of f?}
\end{definition}

\begin{definition}[Path]
    An \textit{infinite path} $\pi$ in a DTMC $\MC$ is a sequence $\pi = s_i \rightarrow s_j \rightarrow \dots$ of states with $s_i,s_j \in S$, such that $\PF(\pi[k],\pi[k+1]) > 0$, where $\pi[k]\in S$ denotes the $k$-th element of the sequence/path for all $k\in \NN_0$. A \textit{finite path} is a finite prefix of an \textit{infinite path}.
\end{definition}
Prob distr. of paths


\begin{definition}[(Step-Bounded) Constrained Reachability (Probabilities?)]
    Let $T,R \subseteq S^{\MC}$ be two sets of states in a DTMC $\MC$. We now define the probability of only visiting states in R until eventually a state in T is visited as $Pr(R\cup T)$, where
    \[R\cup T = \{\pi \in \Path|\exists k\in \NN: \pi[k]\in T\land \forall i <k: \pi [i]\in R\}\]
    Furthermore, the probability of only visint states in $R$ until a state in $T$ is visited at step $k\in \NN$ is given by $Pr(R\cup^{=k})$ \dots (step-bounded defs necessary?)
\end{definition}

\begin{definition}[Reachability Probability]
    The reachability probability of a set $T \subseteq S^{\MC}$ in a DTMC $\MC$ is given by $Pr(\Diamond T)$, where $\Diamond T = S\cup T$
\end{definition}
As $S^{\MC}$ defines all states in $\MC$, this simply defines the probability of visiting a state $t\in T$ at an arbitrary step in the future, i.e. if $Pr(\Diamond T) >0$ at least one state $t\in T$ is reachable. \todo{how to define/denote initial state. Unterscheiden, wenn t initial oder nicht}
In case $T = \{t\}$ is a singleton set, we omit the braces and denote $\Diamond t$ instead of $\Diamond \{t\}$ (and $S \cup t$ instead of $S\cup T$)

$\sum_{\Path}$ DEFINIEREN (?)

We use the LTL unary prefix operator $\cCircle$ to denote the temporal modality "next", meaning 
\begin{definition}[Next Step Probability]\label{def:next_step_probability}
    Let $\MC$ be a DTMC and let $E\in {\Path}$ be a (countable) set of infinite paths in $\MC$.\todo{Difference Paths and $\sum_Paths$. What is necessary here, seems to be that both were used in MA} 
    We now define:
    \[\cCircle E = \{\pi[2] \rightarrow \dots \in \Path|\pi[1]\rightarrow \pi[2] \rightarrow \dots \in E\}\]
    The next step probability with respect to $E$ is then defined as 
    \[Pr(\cCircle E) = Pr()\]
    \todo{Pr(OE) muss Pr sein und nicht set von paths oder (falsch in MA?)}
\end{definition}

\begin{definition}[(Bottom) Strongly Connected Component]
    A set of states $C \subseteq S^{\MC}$ of a DTMC $\MC$ is strongly connected if and only if for every pair of states $s,t\in S$, there exists a path $\pi \in \Path$, such that $\pi[0]=s$ and $\pi[k]$ for some $k\in \NN$. \todo{use induced graph notation additionally}\\
    If $C$ is strongly connected and no proper subset, i.e. $C' \subset C$, is strongly connected, then $C$ is called a strongly connected component (SCC).\\
    Furthermore, an SCC $C$ is called a bottom stronlgy connected component (BSCC) if there exists no $k\in \NN$ and no path $\pi[0] \rightarrow \dots \rightarrow \pi[k] \in \Path$ such that $\pi[0]\in C$ and $\pi[k] \notin C$.
\end{definition}
In other words, once the model enters a BSCC, it cannot leave it again.


SCC (minimal amount of states of set such that still strongly connec)
(Acyclic DTMC)

The states can be classified \dots
\begin{definition}[Transient and recurrent states]\label{def:transient_and_recurrent_states}
    A state $s \in S$ is called recurrent if and only if $Pr_s(\cCircle \Diamond s)=1$.\todo{explain: meaning of notation $Pr_s$} All states that are not recurrent are called transient.
\end{definition}
In other words, a state is recurrent precisely when we are guaranteed to reach it from all of its direct successors, meaning once we reach a recurrent state, we are bound to return to it. (Thus reaching it infinitely often) Using these definitions, we can split the states $S$ of a DTMC into two sets with $S_{tr}$ containing all of the transient states and $S_{re}$ containing all of the recurrent states.

\begin{definition}[Absorbing state]
    A states $s$ of a DTMC $\MC$ is refered to as an absorbing state if and only if $\PF^{\MC}(s,s) = 1$.
\end{definition}
This means, once the model enters an absorbing state, it cannot leave it. Notice, that absorbing states are a special type of recurrent states, as the only direct successor of an absorbing state $s\in S$ is the $s$ itself which trivially has $Pr_s(\Diamond s)=1$.

Using BSCCs we can define a special class of DTMCs.
\begin{definition}[Absorbing DTMC]
    A DTMC $\MC$ is called absorbing if and only if all BSSCs are singleton sets consisting of an absorbing state.    
\end{definition}

\begin{definition}[Reducible DTMC]
    A DTMC $\MC$ is called reducible if and only if 
\end{definition}

\begin{theorem}
    For any DTMC $\MC$, $s\in S$ with $\iota_{init}(s) > 0$ it holds that $Pr(\Diamond S_{re})=1$
\end{theorem}
(Limit Behaviour?)

Canonical Form

(Reward Structure)  

\subsubsection{Expected Visiting Times}
Intuitively, the notion of EVTs describes how often we expect the model to be in a given state. To formally introduce EVTs, we fix a DTMC $\MC$ and start by introducing \textit{visiting times} along paths.
\begin{definition}[Visiting Times of Infinite Paths]
    The function $vt_s:\Path \rightarrow \RR$ describes the \textit{visiting time} of a state $s\in S$ along infinite paths in $\MC$, where for $\pi\in \Path$
    \[vt_s(\pi)=\sum_{k=0}^{\infty}\llbracket \pi[k]=s \rrbracket\]
\end{definition}

Borel Measurability

\begin{definition}[Expected Visiting Times in DTMCs]
    For a state $s\in S$, the \textit{expected visiting time} of $s$ is given by
    \[\Exp^{\MC}[vt_s]=\int_{\pi\in\Path}vt_s(\pi)dPr^{\MC}(\pi)\]   
\end{definition}
OMSSION OF $^{\MC}$

Using the aforementioned classification of the states of $\MC$ into transient and recurrent states in \cref{def:transient_and_recurrent_states}, we can already give an expression for the latter.

\begin{lemma}[EVT of Recurrent State]
    For all recurrent states $s\in S_{re}$ in $\MC$, it holds that
    \[\Exp[vt_s] = 
    \begin{cases*}
        \infty & if s is reachable,\\
        0 & else.
    \end{cases*}\]
\end{lemma}
(Proof??)
Thus, we will focus on finding the non-trivial expected visiting times of the transient states of DTMCs.



EVTs Def.
--> Wie berechnen

\subsubsection{Interval Markov Chains}  
In this thesis, we will focus on Interval Markov Chains. These are defined as follows.
\begin{definition}[Interval Markov Chain]
    An Interval-valued Discrete-time Markov chain (IDTMC) is a tuple $\IMC=(S,\PFupper,\PFlower,\Initupper,\Initlower)$, where
    \begin{enumerate}[label=(\roman*)]
        \item $S$ is a countable, nonempty set of states,
        \item $\PFupper: S \times S \rightarrow [0,1]$ is the transition probability function, where $\PFupper(s,s')$ gives the upper bound on the transition probability from state s to s',
        \item $\PFlower: S \times S \rightarrow [0,1]$ is the transition probability function, where $\PFlower(s,s')$ gives the lower bound on the transition probability from state s to s',
        \item $\Initupper: S \rightarrow [0,1]$ is the upper bound on the initial distribution, and
        \item $\Initlower: S \rightarrow [0,1]$ is the upper bound on the initial distribution.
    \end{enumerate}
\end{definition}
\todo{IDTMC: Interval-valued DTMC}
IMC Def

There are two commonly used semantics, which we will consider in this paper. The first one is the \textit{once-and-for-all}, see \cite{Semantics_Quant-Analysis,Semantics_Model-Checking_omega-Regular,Semantics_LTL_Model_Checking}. 

EVTs für IMCs

\begin{definition}[Once-and-for-all Semantic]
    An arbitrary DTMC $\MC$ satisfies an IDTMC $\IMC$ if and only if $\PFlower^{\IMC}(s,s')\leq\PF^{\MC}(s,s')\leq\PFupper^{\IMC}(s,s')$ and $\Initlower^{\IMC}(s,s')\leq\iota_{init}^{\MC}(s,s')\leq\Initupper^{\IMC}(s,s')$.
\end{definition}
Thus, any DTMC whose initial distribution and transition probabilities lie in the intervals defined by the IDTMC is satisfactory.

\begin{definition}[IMDP semantic]
    An arbitrary DTMC $\MC=(T,\PF,\iota_{init})$ satisfies an IDTMC $\IMC = (S,\PFupper,\PFlower,\Initupper,\Initlower)$ if and only if there exists an injective mapping $\tau: T\rightarrow S$ such that
    \begin{enumerate}[label=(\roman*)]
        \item $\tau(t) = s\in S$ for all $t\in T \implies \Initlower(s) \leq \iota_{init}(t)\leq \Initupper(s)$
        \item $\tau(t) = s\in S$ for all $t\in T \implies \PFlower(s) \leq \PF(t)\leq \PFupper(s)$
    \end{enumerate}
\end{definition}
\dots examples & intuition

(Explain how they describe/define it: the minor differences between them)


\section{Conclusion and Outlook}


% The list of references is provided as `references.bib`
\bibliographystyle{unsrt}
\bibliography{IEEEabrv,references}

\end{document}